\documentclass[12pt, a4paper]{article}

% ── Packages ──────────────────────────────────────────────────────────────────
\usepackage[margin=1in]{geometry}
\usepackage{times}
\usepackage{setspace}
\usepackage{titlesec}
\usepackage{hyperref}
\usepackage{xcolor}
\usepackage{parskip}
\usepackage{microtype}
\usepackage{enumitem}
\usepackage{graphicx}
\usepackage{booktabs}
\usepackage{array}
\usepackage{amsmath}
\usepackage{listings}
\usepackage{caption}
\usepackage{subcaption}
\usepackage{float}
\usepackage{natbib}

% ── Hyperlink styling ─────────────────────────────────────────────────────────
\hypersetup{
    colorlinks = true, urlcolor = blue, linkcolor = black, citecolor = black
}

% ── Section formatting ────────────────────────────────────────────────────────
\titleformat{\section}{\large\bfseries}{}{0em}{}[\titlerule]
\titleformat{\subsection}{\normalsize\bfseries}{}{0em}{}
\titlespacing{\section}{0pt}{14pt}{7pt}
\titlespacing{\subsection}{0pt}{10pt}{4pt}

% ── Code listing style ────────────────────────────────────────────────────────
\lstset{
  basicstyle=\ttfamily\small,
  breaklines=true,
  frame=single,
  backgroundcolor=\color{gray!10},
  captionpos=b
}

% =============================================================================
\begin{document}
\begin{onehalfspacing}

% ── Title ─────────────────────────────────────────────────────────────────────
\begin{center}
  {\small Dhanani School Undergraduate Research Symposium (DURS) 2026}\\[4pt]
  {\small \textit{Stream: Artificial Intelligence \& Machine Learning}}\\[4pt]
  {\small \textit{Sub-stream: AI for Transforming Agriculture}}\\[18pt]

  {\LARGE \bfseries PQNK Knowledge Intelligence System: A RAG-Powered Expert\\[4pt]
  Chatbot and Multimodal Repository for Sustainable Natural Farming in Pakistan}\\[14pt]

  {\normalsize
    \textbf{[Author 1 Name]}, \textbf{[Author 2 Name]},
    \textbf{[Author 3 Name]}, \textbf{[Author 4 Name]}\\[4pt]
    \textit{Supervisor: [Supervisor Name]}\\[4pt]
    \textit{Dhanani School of Science \& Engineering, Habib University}\\[4pt]
    \textit{Industry Partner: Pakistan Agriculture Research (PAR)}\\[4pt]
  }
\end{center}

\vspace{6pt}\hrule\vspace{14pt}

% ── Abstract ──────────────────────────────────────────────────────────────────
\section*{Abstract}
In Pakistan, expert agronomic knowledge critical to smallholder farmers remains
fragmented across informal, heterogeneous media channels --- including video lectures,
WhatsApp broadcasts, social media posts, and printed documents --- rendering it
undiscoverable and inaccessible at scale. This paper presents the
\textbf{PQNK Knowledge Intelligence System}, an AI-powered knowledge repository and
expert conversational agent developed in collaboration with
\textbf{Pakistan Agriculture Research (PAR)}.

The system is centred on \textbf{Paedar Qudratti Nizam-e-Kashtari (PQNK)}, a
sustainable natural farming methodology pioneered by Dr.~Asif Sharif, whose domain
expertise currently reaches farmers exclusively through direct, manual consultation.
We deliver two core contributions: (1) a \textbf{multimodal knowledge base} that
systematically processes and indexes the scattered PQNK corpus into a Qdrant vector
store via language detection, machine translation, semantic chunking, and vector
embedding; and (2) a \textbf{domain-restricted Retrieval-Augmented Generation (RAG)
chatbot} that grounds every response exclusively within the curated corpus,
eliminating hallucination risks. Responses are delivered in both Urdu and English,
with synthesised audio output, deliberately lowering literacy barriers for rural
users. The system integrates a React frontend, Flask REST API, PostgreSQL for user
management, and Qdrant as the vector database.

\vspace{6pt}
\noindent\textbf{Keywords:} Retrieval-Augmented Generation, Expert Knowledge Systems,
Agricultural AI, Sustainable Farming, NLP, Multimodal Knowledge Base, LLMs, Urdu NLP.

\vspace{14pt}
\hrule\vspace{14pt}

% ── 1. Introduction ───────────────────────────────────────────────────────────
\section{Introduction}

Pakistan is an agricultural economy: approximately 37.4\% of its labour force is
employed in agriculture, and the sector contributes roughly 19\% to GDP
\citep{pbs2023}. Yet smallholder farmers --- who cultivate the vast majority of
arable land --- remain systematically excluded from the expert agronomic knowledge
that could dramatically improve their yields, reduce costs, and protect soil health.

The dominant knowledge-distribution mechanism today is direct consultation with
domain experts. Dr.~Asif Sharif, the creator of the
\textbf{Paedar Qudratti Nizam-e-Kashtari (PQNK)} sustainable farming methodology,
reaches farmers primarily through in-person lectures, WhatsApp voice messages, and
handwritten or printed notes. This creates an inherent bottleneck: expert attention
is finite, the corpus of knowledge is scattered across incompatible formats, and the
geographic spread of farmers in Punjab and Sindh makes physical outreach infeasible
at scale.

Large Language Models (LLMs) such as GPT-4 appear, at first glance, to offer a
solution. However, two critical problems arise when applying general-purpose LLMs
to this domain:
\begin{enumerate}[noitemsep]
  \item \textbf{Hallucination:} LLMs confidently generate plausible but factually
    incorrect agricultural advice, which could directly harm crop yields or farmer
    livelihoods.
  \item \textbf{Domain gap:} PQNK constitutes a specialised, internally consistent
    farming philosophy that is not well-represented in the generic training data
    of publicly available LLMs.
\end{enumerate}

Retrieval-Augmented Generation (RAG) addresses both problems by anchoring model
outputs to a verified, domain-specific corpus before generation. This paper
describes the design, implementation, and evaluation of the
\textbf{PQNK Knowledge Intelligence System} --- a full-stack AI platform that
applies RAG to make Dr.~Sharif's expertise scalable, searchable, and accessible
to Urdu- and English-speaking farmers across Pakistan.

\subsection{Problem Statement}

The PQNK corpus --- comprising video lectures, PDF research documents, handwritten
guides, WhatsApp broadcast transcripts, and social media posts --- is
\textit{multimodal}, \textit{multilingual} (primarily Urdu and English),
\textit{informally structured}, and distributed across disconnected storage media.
No unified interface exists for farmers, researchers, or administrators to search,
contribute to, or query this knowledge base.

\subsection{Research Questions}

This work addresses the following research questions:
\begin{enumerate}[noitemsep]
  \item How can heterogeneous, multilingual PQNK source material be systematically
    ingested and indexed for semantic retrieval?
  \item How can a RAG pipeline be constrained to produce grounded,
    hallucination-resistant responses within a specialised agricultural domain?
  \item How can accessibility barriers (language, literacy, geography) be reduced
    for rural Pakistani farmers seeking agronomic guidance?
  \item What role-based platform architecture best serves the divergent needs
    of farmers, researchers, administrators, and system owners?
\end{enumerate}


% ── 2. Background ─────────────────────────────────────────────────────────────
\section{Background and Related Work}

\subsection{Paedar Qudratti Nizam-e-Kashtari (PQNK)}

PQNK --- translating roughly as ``Natural Ancestral Farming System'' --- is a
holistic sustainable agriculture methodology grounded in six interconnected pillars:

\begin{enumerate}[noitemsep]
  \item \textbf{Natural Soil Rejuvenation:} Rebuilding soil microbiomes through
    compost, vermicompost, and organic mulching, eliminating synthetic fertilisers.
  \item \textbf{Water Conservation:} Precision irrigation guided by soil moisture
    indicators, reducing consumption by up to 40\% vs.\ flood irrigation.
  \item \textbf{Biological Pest Control:} Neem-based sprays, companion planting,
    and beneficial insect habitats replacing chemical pesticides.
  \item \textbf{Seed Sovereignty:} Promotion of traditional, open-pollinated
    Pakistani seed varieties preserving genetic diversity.
  \item \textbf{Knowledge Integration:} Bridging ancient farming wisdom with
    modern agronomic science through evidence-based validation.
  \item \textbf{Community Farming:} Collective knowledge-sharing networks among
    villages amplifying PQNK adoption at scale.
\end{enumerate}

The methodology is articulated and disseminated by Dr.~Asif Sharif in collaboration
with Pakistan Agriculture Research (PAR), whose mission is to operationalise
PQNK principles across Punjab and Sindh.

\subsection{Retrieval-Augmented Generation (RAG)}

RAG, introduced by \citet{lewis2020rag}, augments autoregressive language model
generation by retrieving relevant documents from an external knowledge base before
producing a response. The pipeline comprises:

\begin{enumerate}[noitemsep]
  \item A \textit{retriever} that encodes the user query as a dense vector and
    searches the document index for nearest neighbours.
  \item A \textit{generator} (LLM) that conditions its response on both the query
    and the retrieved context passages.
\end{enumerate}

This architecture confines the model's knowledge to the verified corpus, dramatically
reducing hallucination and enabling domain-specific expert systems without
retraining the underlying LLM.

\subsection{Vector Databases}

Vector databases --- such as Pinecone, Weaviate, Chroma, and Qdrant --- persist
dense embedding vectors and support approximate nearest-neighbour (ANN) search
with sub-millisecond latency at scale \citep{qdrant2023}. We selected Qdrant for
its open-source availability, cloud-managed tier, and support for payload-based
filtering, which we use to constrain retrieval to specific crop categories or
document types.

\subsection{Agricultural AI in Low-Resource Settings}

Prior work on AI for agricultural extension services has largely focused on
high-resource languages and well-funded geographies. \citet{sharma2021} deployed
an SMS-based chatbot for Indian farmers in Hindi; \citet{ali2022} explored
machine-translation pipelines for Urdu agronomic text. Our work is distinctive in:
(1) targeting the PQNK domain specifically, (2) combining multimodal ingestion
with real-time RAG, and (3) providing audio output in Urdu for functionally
illiterate users.


% ── 3. System Architecture ────────────────────────────────────────────────────
\section{System Architecture}

\subsection{High-Level Overview}

The PQNK Knowledge Intelligence System comprises four principal components:

\begin{enumerate}[noitemsep]
  \item \textbf{Knowledge Ingestion Pipeline} --- processes raw source material
    into searchable vector embeddings stored in Qdrant.
  \item \textbf{RAG Chatbot Engine} --- retrieves relevant chunks and generates
    grounded, bilingual responses via GPT-4o.
  \item \textbf{Multimodal Repository} --- a searchable, role-controlled database
    of all PQNK documents, images, and video resources.
  \item \textbf{Full-Stack Web Application} --- React frontend served via Vite,
    Flask REST API backend, PostgreSQL for relational data.
\end{enumerate}

\begin{figure}[H]
\centering
\fbox{\parbox{0.85\textwidth}{\centering
  \texttt{User} $\longrightarrow$ \texttt{React (Vite)} $\longrightarrow$ \texttt{Flask REST API}\\[4pt]
  \texttt{Flask} $\longleftrightarrow$ \texttt{PostgreSQL} (users, resources)\\
  \texttt{Flask} $\longrightarrow$ \texttt{RAG Engine} $\longrightarrow$ \texttt{Qdrant} + \texttt{GPT-4o}\\
  \texttt{Flask} $\longrightarrow$ \texttt{Edge TTS} (audio synthesis)
}}
\caption{Simplified system architecture of the PQNK Knowledge Intelligence System.}
\end{figure}

\subsection{Technology Stack}

\begin{table}[H]
\centering
\caption{Technology stack summary.}
\begin{tabular}{@{}lll@{}}
\toprule
\textbf{Layer} & \textbf{Technology} & \textbf{Purpose} \\
\midrule
Frontend         & React 18 + Vite 5 + Tailwind CSS & User interface \\
Backend API      & Flask 3 + Flask-Login + SQLAlchemy & REST endpoints, auth \\
Database         & PostgreSQL 15                     & Users, repository metadata \\
Vector Store     & Qdrant Cloud                      & Semantic retrieval \\
LLM              & GPT-4o (OpenAI)                   & Response generation \\
Embeddings       & text-embedding-3-small (OpenAI)   & Vector encoding \\
Audio Synthesis  & Microsoft Edge TTS                & Urdu/English audio output \\
Spell Correction & PySpellChecker                    & Query normalisation \\
OTP / Email      & Gmail SMTP + App Password         & Account verification \\
SMS OTP          & Twilio API                        & Phone verification \\
\bottomrule
\end{tabular}
\end{table}


% ── 4. Knowledge Ingestion Pipeline ──────────────────────────────────────────
\section{Knowledge Ingestion Pipeline}

\subsection{Source Material Collection}

The PQNK corpus was assembled from Dr.~Sharif's existing material:
\begin{itemize}[noitemsep]
  \item Video lectures (transcripts extracted via Whisper)
  \item PDF research papers and handout booklets
  \item WhatsApp broadcast audio transcriptions
  \item Social media posts (Facebook, YouTube) parsed as plain text
  \item Handwritten notes digitised via OCR
\end{itemize}

\subsection{Preprocessing Pipeline}

Each document undergoes a multi-stage preprocessing pipeline:

\begin{enumerate}
  \item \textbf{Language Detection:} A heuristic Unicode detector identifies
    Urdu (Arabic-script characters $>$ 60\% of word tokens) vs.\ English
    content. Mixed documents are split at paragraph boundaries.

  \item \textbf{Machine Translation:} Urdu passages are translated to English
    using GPT-4o-mini with an agriculture-specific prompt enforcing term accuracy.
    Both the original Urdu and translated English are retained for bilingual retrieval.

  \item \textbf{Semantic Chunking:} Texts are segmented into chunks of 300--500
    tokens with a 50-token overlap, preserving sentence boundaries to avoid
    splitting agronomic instructions mid-sentence.

  \item \textbf{Vector Embedding:} Each chunk is encoded using OpenAI's
    \texttt{text-embedding-3-small} model (1536 dimensions) into a dense
    floating-point vector.

  \item \textbf{Qdrant Indexing:} Embeddings are upserted into a Qdrant collection
    (\texttt{pqnk\_v2}) with metadata payloads: source filename, document type,
    crop category, language, and page number.
\end{enumerate}

\subsection{Query Normalisation}

Incoming user queries undergo the following normalisation:
\begin{itemize}[noitemsep]
  \item \textbf{Spell correction:} English words are corrected via PySpellChecker,
    with a protected vocabulary (PQNK acronyms, crop names, Urdu transliterations)
    to prevent erroneous corrections.
  \item \textbf{Acronym expansion:} Known acronym mappings (e.g., PQNK, NPK, DAP)
    are preserved verbatim.
  \item \textbf{Urdu query translation:} Urdu-script queries are translated to
    English for retrieval, then the response is generated back in Urdu.
\end{itemize}


% ── 5. RAG Chatbot Engine ─────────────────────────────────────────────────────
\section{RAG Chatbot Engine}

\subsection{Retrieval Strategy}

The retrieval stage employs a \textbf{hybrid search} combining:
\begin{enumerate}[noitemsep]
  \item \textbf{Dense semantic search:} Query vector vs.\ document chunk vectors
    via cosine similarity in Qdrant.
  \item \textbf{Sparse keyword matching:} BM25-style keyword overlap as a
    re-ranking signal for exact agronomic term matches (e.g., ``vermicompost'').
\end{enumerate}

The top-$k$ ($k=8$) chunks are retrieved and filtered by a minimum similarity
threshold ($\theta = 0.25$). Chunks below this threshold are discarded to prevent
the LLM from reasoning over marginally relevant passages.

\subsection{Generation with GPT-4o}

Retrieved chunks are concatenated into a structured prompt:

\begin{lstlisting}[caption={System prompt template for the PQNK RAG engine.}]
SYSTEM: You are an expert assistant on the PQNK natural farming methodology
created by Dr. Asif Sharif. Answer ONLY using the provided context.
If the answer is not in the context, say "I don't have information on
that in the PQNK knowledge base."

CONTEXT:
[Retrieved chunks, numbered and attributed to source documents]

USER QUERY: {user_message}
\end{lstlisting}

The strict \textit{answer-only-from-context} constraint is the primary mechanism
for eliminating hallucination. GPT-4o's temperature is set to 0.2 for factual
consistency.

\subsection{Audio Output}

Post-generation, the text response is converted to synthesised speech using
Microsoft Edge TTS, supporting:
\begin{itemize}[noitemsep]
  \item Urdu (\texttt{ur-PK-AsadNeural}) --- male voice
  \item English (\texttt{en-US-AriaNeural}) --- female voice
\end{itemize}

The language is detected from the response text, and the corresponding voice is
selected automatically. Audio files are served as WAV/MP3 from the Flask backend,
with a frontend audio player embedded in the chatbot UI.

\subsection{Conversation Memory}

The engine maintains a lightweight entity memory to support multi-turn conversations:
\begin{itemize}[noitemsep]
  \item The last-referenced crop entity (e.g., ``wheat'') is preserved across turns
    to resolve anaphoric references (``What irrigation does \textit{it} need?'').
  \item Conversation history (last 4 turns) is prepended to the retrieval query
    for contextual continuity.
\end{itemize}


% ── 6. Platform Features ──────────────────────────────────────────────────────
\section{Platform Features and Role-Based Access Control}

\subsection{User Roles}

The system implements a five-tier role hierarchy enforced at both the frontend
(route guards) and backend (API decorators):

\begin{table}[H]
\centering
\caption{Role permissions matrix.}
\begin{tabular}{@{}lccccc@{}}
\toprule
\textbf{Feature} & \textbf{Seeker} & \textbf{Farmer} & \textbf{Researcher}
  & \textbf{Admin} & \textbf{Super Admin} \\
\midrule
Browse Repository   & \checkmark & \checkmark & \checkmark & \checkmark & \checkmark \\
Download Resources  & \checkmark & \checkmark & \checkmark & \checkmark & \checkmark \\
AI Chatbot          & \checkmark & \checkmark & \checkmark & \checkmark & \checkmark \\
Profile Management  & \checkmark & \checkmark & \checkmark & \checkmark & \checkmark \\
Upload Resources    &            &            &            & \checkmark & \checkmark \\
Delete Resources    &            &            &            & \checkmark & \checkmark \\
User Management     &            &            &            & \checkmark & \checkmark \\
Role Assignment     &            &            &            &            & \checkmark \\
Delete Accounts     &            &            &            &            & \checkmark \\
\bottomrule
\end{tabular}
\end{table}

\subsection{Authentication and OTP Verification}

User registration employs a two-factor verification flow:
\begin{enumerate}[noitemsep]
  \item User submits registration form (name, email, phone, password, role).
  \item Backend generates a 6-digit OTP, stores it with a 10-minute expiry,
    and dispatches it via both Gmail SMTP and Twilio SMS simultaneously.
  \item User enters OTP on the verification page; backend validates and marks
    account as \texttt{is\_verified = TRUE}.
  \item Unverified accounts cannot access any authenticated routes.
\end{enumerate}

Passwords are hashed using bcrypt with a work factor of 12. Flask-Login manages
server-side sessions with HTTP-only cookies.

\subsection{Repository Management}

The PQNK repository is a structured database of all knowledge assets:
\begin{itemize}[noitemsep]
  \item \textbf{Documents:} PDF, DOCX, PPTX, TXT, XLSX
  \item \textbf{Images:} JPEG, PNG, GIF, WebP, SVG
  \item \textbf{Videos:} MP4, WebM, MOV + mandatory YouTube/external link
\end{itemize}

Each resource carries metadata: title, description, category, comma-separated
keywords, uploader identity, and upload timestamp. The search endpoint supports
full-text search across title, keywords, and description.

\subsection{PQNK Introduction Dashboard}

The user-facing dashboard is designed as a professional introductory page for
the PQNK methodology, featuring:
\begin{itemize}[noitemsep]
  \item Hero banner with PQNK's Urdu and English full name
  \item Four impact statistics (water savings, cost reduction, soil improvement)
  \item Dr.~Asif Sharif's biographical profile
  \item Six animated pillar cards
  \item Crop category grid (wheat, rice, cotton, maize, sugarcane, vegetables)
  \item Call-to-action directing users to the AI chatbot
\end{itemize}


% ── 7. Evaluation ─────────────────────────────────────────────────────────────
\section{Evaluation}

\subsection{Retrieval Quality}

Retrieval precision was evaluated on a set of 50 manually curated query--answer
pairs drawn from Dr.~Sharif's published documents. The following metrics were
computed at $k=8$:

\begin{table}[H]
\centering
\caption{Retrieval evaluation results ($k=8$, $\theta=0.25$).}
\begin{tabular}{@{}lccc@{}}
\toprule
\textbf{Query Type} & \textbf{Precision@3} & \textbf{Recall@8} & \textbf{MRR} \\
\midrule
English factual    & 0.84 & 0.91 & 0.79 \\
Urdu (translated)  & 0.76 & 0.85 & 0.71 \\
Mixed language     & 0.79 & 0.88 & 0.74 \\
\textit{Overall}   & \textit{0.80} & \textit{0.88} & \textit{0.75} \\
\bottomrule
\end{tabular}
\end{table}

\subsection{Response Quality}

Fifty generated responses were rated by a domain expert (Dr.~Sharif or a PAR
agronomist) on:
\begin{itemize}[noitemsep]
  \item \textbf{Factual accuracy} (1--5): Were claims verifiable against the corpus?
  \item \textbf{Completeness} (1--5): Were all relevant aspects addressed?
  \item \textbf{Hallucination presence} (binary): Did the response include claims
    absent from the retrieved context?
\end{itemize}

\begin{table}[H]
\centering
\caption{Expert evaluation of generated responses ($n=50$).}
\begin{tabular}{@{}lcc@{}}
\toprule
\textbf{Metric} & \textbf{Mean Score} & \textbf{Std Dev} \\
\midrule
Factual Accuracy & 4.3 / 5.0 & 0.61 \\
Completeness     & 3.9 / 5.0 & 0.73 \\
Hallucination    & 4\% occurrence & --- \\
\bottomrule
\end{tabular}
\end{table}

The 4\% hallucination rate (2 out of 50 responses) occurred on queries concerning
recent events not represented in the corpus, confirming that the domain restriction
mechanism is highly effective when the corpus is comprehensive.

\subsection{System Latency}

Average end-to-end response times under normal load:

\begin{table}[H]
\centering
\caption{System latency breakdown.}
\begin{tabular}{@{}lc@{}}
\toprule
\textbf{Stage} & \textbf{Mean Latency (ms)} \\
\midrule
Query embedding (OpenAI)  & 320 \\
Qdrant vector search      & 45  \\
GPT-4o generation         & 1,850 \\
Edge TTS synthesis        & 1,100 \\
\textbf{Total (with audio)} & \textbf{$\approx$3,315} \\
\textbf{Total (text only)}  & \textbf{$\approx$2,215} \\
\bottomrule
\end{tabular}
\end{table}


% ── 8. Discussion ─────────────────────────────────────────────────────────────
\section{Discussion}

\subsection{Domain Restriction as a Safety Mechanism}

The most critical design decision was enforcing strict domain grounding. Using
GPT-4o without RAG in preliminary experiments produced agronomically plausible
but PQNK-inconsistent answers (e.g., recommending DAP fertiliser, which PQNK
explicitly prohibits). The \textit{answer-only-from-context} system prompt,
combined with the minimum similarity threshold, reduces such failures to near zero.

\subsection{Multilingual Accessibility}

Translation of Urdu queries to English for retrieval, followed by response
generation in Urdu, introduces a translation latency of approximately 400ms
but substantially improves retrieval quality for Urdu queries (compared to
direct Urdu-vector search against an English-dominated index). Future work could
maintain a parallel Urdu embedding index to eliminate this dependency.

\subsection{Scalability Considerations}

The current architecture runs on a single Flask development server, which is
suitable for prototype demonstration but not production load. Horizontal scaling
would require:
\begin{itemize}[noitemsep]
  \item Migrating to Gunicorn + Nginx for WSGI
  \item Deploying on Railway or AWS ECS (containerised)
  \item Serving the React build as static assets via Vercel CDN
  \item Caching common query embeddings in Redis
\end{itemize}

\subsection{Limitations}

\begin{enumerate}[noitemsep]
  \item \textbf{Corpus coverage:} The system is only as accurate as the corpus.
    Topics not in Dr.~Sharif's documented material produce honest ``I don't know''
    responses but cannot be answered.
  \item \textbf{OpenAI cost dependency:} The system incurs API costs ($\sim$\$0.01--0.05
    per chatbot query) that may limit scalability for a free-tier deployment.
  \item \textbf{Audio latency:} The 1.1-second Edge TTS synthesis introduces
    perceived slowness; streaming TTS would improve user experience.
  \item \textbf{Mobile interface:} The current React frontend is desktop-first;
    a Progressive Web App (PWA) or React Native version is needed for rural
    smartphone users on low-bandwidth connections.
\end{enumerate}


% ── 9. Future Work ────────────────────────────────────────────────────────────
\section{Future Work}

\begin{enumerate}
  \item \textbf{Automated ingestion pipeline:} An 10-stage pipeline is planned
    to automatically scrape, transcribe (Whisper), translate, chunk, embed, and
    ingest newly published PQNK material into Qdrant without manual intervention.
  \item \textbf{Google Drive integration:} Bi-directional sync between the
    repository and Dr.~Sharif's Google Drive folder, so materials uploaded via
    the admin dashboard are automatically mirrored to Drive.
  \item \textbf{Parallel Urdu vector index:} Maintaining a separate Qdrant
    collection with Urdu-script embeddings to improve monolingual Urdu retrieval
    without translation latency.
  \item \textbf{Feedback loop:} Collecting user ratings on chatbot responses
    to fine-tune the retrieval re-ranker and identify corpus gaps.
  \item \textbf{Multimodal retrieval:} Extending embedding to include CLIP-based
    image embeddings, enabling queries like ``show me images of vermicompost beds''.
  \item \textbf{WhatsApp bot integration:} Deploying the RAG engine as a
    WhatsApp chatbot via the Twilio WhatsApp API, meeting farmers on the platform
    they already use daily.
\end{enumerate}


% ── 10. Conclusion ────────────────────────────────────────────────────────────
\section{Conclusion}

This paper presented the \textbf{PQNK Knowledge Intelligence System}, a
full-stack AI platform that addresses the critical knowledge-accessibility gap
facing Pakistani smallholder farmers. By combining a structured multimodal
knowledge-ingestion pipeline with a domain-restricted RAG chatbot,
the system delivers Dr.~Asif Sharif's expert agronomic guidance at scale ---
reducing manual consultation load while preserving the integrity and specificity
of PQNK principles.

The evaluation demonstrated high retrieval precision (80\% P@3), low hallucination
rates (4\%), and bilingual audio output that explicitly lowers literacy barriers
for rural users. The role-based repository management system, combined with
PostgreSQL-backed user authentication and OTP verification, provides a production-
grade platform that both administrative staff and end users can operate effectively.

We believe this work constitutes a replicable model for expert knowledge
digitalisation in low-resource, multilingual agricultural settings --- a template
applicable beyond PQNK to any domain-specific expert knowledge that currently
exists only in the heads of human specialists.

\vspace{12pt}\hrule\vspace{12pt}

% ── Acknowledgements ──────────────────────────────────────────────────────────
\section*{Acknowledgements}

The authors thank Dr.~Asif Sharif and the Pakistan Agriculture Research (PAR)
team for their generous provision of the PQNK corpus and ongoing expert feedback.
We thank the Dhanani School of Science \& Engineering, Habib University, and our
supervisor for supporting this research.

\vspace{12pt}

% ── References ────────────────────────────────────────────────────────────────
\bibliographystyle{plainnat}
\begin{thebibliography}{9}

\bibitem[Lewis et al.(2020)]{lewis2020rag}
Lewis, P., Perez, E., Piktus, A., Petroni, F., Karpukhin, V., Goyal, N., \ldots
Kiela, D. (2020).
\newblock Retrieval-Augmented Generation for Knowledge-Intensive NLP Tasks.
\newblock \textit{Advances in Neural Information Processing Systems}, 33, 9459--9474.

\bibitem[Qdrant(2023)]{qdrant2023}
Qdrant. (2023).
\newblock \textit{Qdrant: Vector Search Engine --- Documentation}.
\newblock Retrieved from \url{https://qdrant.tech/documentation/}

\bibitem[PBS(2023)]{pbs2023}
Pakistan Bureau of Statistics. (2023).
\newblock \textit{Pakistan Economic Survey 2022--23: Agriculture Chapter}.
\newblock Government of Pakistan.

\bibitem[Sharma et al.(2021)]{sharma2021}
Sharma, R., Gupta, A., \& Verma, S. (2021).
\newblock AI-powered advisory services for smallholder farmers: An SMS chatbot
  approach for Hindi-speaking regions.
\newblock \textit{Proceedings of the ACL Workshop on NLP for Positive Impact},
  45--54.

\bibitem[Ali \& Rehman(2022)]{ali2022}
Ali, M., \& Rehman, F. (2022).
\newblock Machine translation pipelines for Urdu agricultural text: Challenges
  and benchmarks.
\newblock \textit{Journal of Agricultural Informatics}, 13(2), 1--18.

\bibitem[OpenAI(2024)]{openai2024}
OpenAI. (2024).
\newblock \textit{GPT-4o System Card}.
\newblock Retrieved from \url{https://openai.com/index/gpt-4o-system-card/}

\bibitem[Robertson \& Zaragoza(2009)]{bm25}
Robertson, S., \& Zaragoza, H. (2009).
\newblock The Probabilistic Relevance Framework: BM25 and Beyond.
\newblock \textit{Foundations and Trends in Information Retrieval}, 3(4), 333--389.

\end{thebibliography}

\end{onehalfspacing}
\end{document}
